% ============================================================
\chapter{Matrices and Linear Maps}
\label{ch:matrices-linear-maps}
% ============================================================

\section*{Chapter Preview}
\addcontentsline{toc}{section}{Chapter Preview}

Matrices are not just rectangular tables of numbers. They represent linear
maps once bases have been chosen. This chapter connects the algebra of matrix
multiplication with the geometry of transformations.

\section{Matrix-Vector Multiplication}

Let $A\in\R^{m\times n}$ have columns $a_1,\ldots,a_n$. For
$x=(x_1,\ldots,x_n)^T$, define
\[
  Ax=x_1a_1+\cdots+x_na_n.
\]
Thus $Ax$ is a linear combination of the columns of $A$.

\begin{example}
If
\[
  A=\mat{1&2\\3&4\\0&1},\qquad x=\mat{5\\-1},
\]
then
\[
  Ax=5\mat{1\\3\\0}-1\mat{2\\4\\1}
    =\mat{3\\11\\-1}.
\]
\end{example}

\section{Linear Maps}

\begin{definition}
  A map $T:V\to W$ between vector spaces over the same field is \emph{linear}
  if
  \[
    T(u+v)=T(u)+T(v),
    \qquad
    T(\lambda v)=\lambda T(v)
  \]
  for all $u,v\in V$ and scalars $\lambda$.
\end{definition}

Equivalently, $T(\lambda u+\mu v)=\lambda T(u)+\mu T(v)$ for all choices of
$u,v,\lambda,\mu$.

\begin{proposition}
  If $A$ is an $m\times n$ matrix, then $T(x)=Ax$ is a linear map
  $\R^n\to\R^m$.
\end{proposition}

\begin{proof}
  Matrix-vector multiplication distributes over vector addition and scalar
  multiplication:
  \[
    A(u+v)=Au+Av,\qquad A(\lambda u)=\lambda Au.
  \]
\end{proof}

\section{A Linear Map Is Determined by a Basis}

If $b_1,\ldots,b_n$ is a basis of $V$, then a linear map $T:V\to W$ is
completely determined by the vectors $T(b_1),\ldots,T(b_n)$. Indeed, if
\[
  v=c_1b_1+\cdots+c_nb_n,
\]
then linearity forces
\[
  T(v)=c_1T(b_1)+\cdots+c_nT(b_n).
\]

This is why matrices work: the columns tell us where the basis vectors go.

\section{Matrix Multiplication as Composition}

If $A:\R^n\to\R^m$ and $B:\R^m\to\R^p$, then applying $A$ and then $B$ gives
the composite map
\[
  x\mapsto B(Ax).
\]
The matrix of this composite map is $BA$.

\begin{theorem}
  Matrix multiplication represents composition of linear maps.
\end{theorem}

\begin{proof}
  The $j$-th column of the matrix of the composite is the image of the $j$-th
  standard basis vector:
  \[
    B(Ae_j).
  \]
  Since $Ae_j$ is the $j$-th column of $A$, the columns are exactly the
  columns produced by the usual definition of $BA$.
\end{proof}

\section{Python: Building a Matrix from a Rule}

\begin{lstlisting}
import numpy as np

def matrix_of_T(T, n):
    cols = []
    for j in range(n):
        e = np.zeros(n)
        e[j] = 1.0
        cols.append(T(e))
    return np.column_stack(cols)

def reflect_across_y_equals_x(v):
    return np.array([v[1], v[0]])

A = matrix_of_T(reflect_across_y_equals_x, 2)
print(A)
\end{lstlisting}

\section*{Chapter Summary}
\addcontentsline{toc}{section}{Chapter Summary}

\begin{itemize}
  \item $Ax$ is a linear combination of the columns of $A$.
  \item Linear maps preserve addition and scalar multiplication.
  \item A linear map is determined by its values on a basis.
  \item Matrix multiplication represents composition.
\end{itemize}

\section*{Where This Appears Later}
\addcontentsline{toc}{section}{Where This Appears Later}

Image, kernel, rank, change of basis, eigenvalues, and SVD all build on the
idea that a matrix is a coordinate representative of a linear map.

\begin{labbox}
Lab 2 builds matrices for differentiation, integration, rotations,
reflections, projections, and composition.
\end{labbox}

\input{book/exercises/ch05-exercises}
